% ===================================================================
%  TABEL Nr. 11 — Die stad en sy gedenktekens. — Die brand.
%  Meme regime que les precedents. Premiere de la TROISIEME SERIE :
%  l'apparat porte « DERDE REEKS », comme au tableau 7 « TWEEDE
%  REEKS ».
%
%  DEUX BLOCS PORTENT DEUX ALINEAS SOUS UNE SEULE CLE, t11-c1-12-1 et
%  t11-c2-03-1 : le fac-simile ido y a une coupure de page, pas un
%  paragraphe. Les traductions les soudent et ordo.py connait le cas.
%
%  LE 9 EST LE PLUS GRAS DU LIVRET ET LE PLUS ETRANGE : le
%  compositeur y met en gras des mots-outils entiers — « quale »,
%  « la sistemo di traciono di la », « La », « tirata », « urbo al
%  altra ». Aucun n'est un terme de planche. On les reproduit parce
%  que toutes les autres colonnes les reproduisent.
%
%  ET LE 6 DE LA DEUXIEME SCENE DONNE DEUX FOIS LE RENVOI (81), aux
%  musiciens puis au violon. La planche numerote ainsi ; on ne
%  corrige pas.
%
%  UN MOT SUR LA POLICE DE CE TABLEAU : « jendarmi » (63) et
%  « policisti » (64) sont deux corps distincts, et « kavalkanta
%  jendarmi urbala » (61) un troisieme. L'afrikaans a de quoi les
%  separer — gendarmes, polisiemanne, berede stadspolisie — et il
%  faut le faire, sinon trois renvois de la planche tombent sur le
%  meme mot.
% ===================================================================

%%K t11-apar-1 apar
\VUsaut{2.0mm}
\VUcentre{13.2pt}{90}{DERDE REEKS}
\VUsaut{3.0mm}
\VUfilet{16mm}
\VUsaut{5.6mm}
\VUtitre{15.0pt}{60}{TABEL Nr. 11}
\VUsaut{1.4mm}
\VUpk{11.4pt}{40}{Die stad en sy gedenktekens. --- Die brand.}
\VUsaut{2.2mm}

%%K t11-tit-1 sub
\VUpk{10.2pt}{40}{Die voorstad.}
\VUsaut{3.0mm}

%%K t11-c1-01-1 p
1. --- Ek woon in 'n groot stad wat 100 kilometer van die Oseaan af lê, en wat tog 'n
\VUgras{hawe} \textsuperscript{(1)} van die eerste rang is. Die rivier wat dit omsoom het
'n breedte van 400 meter en vorm 'n baie mooi ree.

%%K t11-c1-02-1 p
2. --- Die hele deel van die stad wat aan die \VUgras{kaaie} \textsuperscript{(2)} lê,
lyk heeltemal maritiem en nywerheidsagtig, en vorm 'n \VUgras{voorstad}
\textsuperscript{(3)} wat net deur seelui en werkers bewoon word. Dié werk in die
\VUgras{fabrieke} \textsuperscript{(4)} en in die \VUgras{gasfabriek}
\textsuperscript{(5)} waarvan die hoë \VUgras{skoorsteen} \textsuperscript{(6)} en die
\VUgras{gashouer} van baie ver af gesien word.

%%K t11-c1-03-1 p
3. --- In die middeleeue was die stad versterk, en 'n mens vind nog enkele oorblyfsels
van sy walle, veral die \VUgras{poort} \textsuperscript{(8)} van die ou \VUgras{versterkte
kasteel} \textsuperscript{(7)} geflankeer deur sy twee \VUgras{torings}
\textsuperscript{(9)} wat nou as \VUgras{gevangenis} gebruik word.

%%K t11-c1-04-1 p
4. --- In daardie hele \VUgras{wyk,} wat baie oud lyk, is net een gebou betreklik modern:
dit is die \VUgras{doeanekantoor} \textsuperscript{(10)}. Dit lê tussen die kaai en die
\VUgras{straatjie} \textsuperscript{(11)} wat na die ou \VUgras{stadhuis}
\textsuperscript{(12)} lei. 'n Mens herken daardie laaste gedenkteken maklik aan die
reusagtige \VUgras{kloktoring} \textsuperscript{(13)} wat die ou stad oorheers.

%%K t11-c1-05-1 p
5. --- Aan die ander kant van die straat \VUgras{staan} 'n gebou wat in dieselfde styl as
die doeanekantoor gebou is: dit is die \VUgras{Beurs} \textsuperscript{(14)}. Een van sy
gewels kyk uit op 'n klein plein waar die \VUgras{prefektuurgebou} \textsuperscript{(15)}
is. Ons stad is naamlik, sonder om 'n hoofstad te wees, tog 'n belangrike
departementshoofstad.

%%K t11-tit-2 sub
\VUsaut{5.0mm}
\VUpk{10.2pt}{40}{Die hoogste deel van die stad.}
\VUsaut{3.4mm}

%%K t11-c1-06-1 p
6. --- Die garnisoen, redelik talryk, is in ruim, baie gesonde \VUgras{kaserne}
\textsuperscript{(16)} gehuisves; hulle strek tot by die hek van die \VUgras{gemeentepark}
\textsuperscript{(17)}. Die \VUgras{boulevard} \textsuperscript{(19)} wat na daardie park
lei loop agter die \VUgras{spaarbank} \textsuperscript{(18)} verby en eindig op die plein
van die \VUgras{katedraal} \textsuperscript{(20)}.

%%K t11-c1-07-1 p
7. --- Al daardie gedenktekens strek hulle soos 'n amfiteater uit op 'n heuweltjie waar
ook ons ruim \VUgras{liseum} \textsuperscript{(21)} staan.

%%K t11-c1-07-2 p
As 'n mens langs 'n effens hellende pad afdaal wat by die Groot straat aansluit, kom 'n
mens by die \VUgras{Geregshof} \textsuperscript{(22)} waarvoor 'n aangename
\VUgras{openbare tuin} \textsuperscript{(26)} hom uitstrek. Daardie \VUgras{plantsoen} is
nie baie groot nie, maar dit maak in die middel van die stad 'n oase van groenigheid en
koelte. Daarom word dit baie deur die stedelinge besoek wat graag onder die tentdak van
die \VUgras{kafee} \textsuperscript{(25)} kom sit, om te luister na die soldate-musikante
wat Donderdae en Sondae in die \VUgras{kiosk} \textsuperscript{(27)} speel wat tussen die
grasperke en die struikgroepe staan.

%%K t11-c1-08-1 p
8. --- Op een van daardie grasperke staan 'n brons \VUgras{standbeeld}
\textsuperscript{(28)}, 'n gedenkteken opgerig ter nagedagtenis van een van ons
stadsgenote wat groot dienste aan sy geboortestad bewys het.

%%K t11-tit-3 sub
\VUsaut{4.0mm}
\VUpk{10.2pt}{40}{Die vervoermiddels.}
\VUsaut{3.4mm}

%%K t11-c1-09-1 p
9. --- Tot nou toe word ons stad nog nie met elektriese lig verlig nie, dit het net
\VUgras{gaslampe} \textsuperscript{(29)}. Maar die \VUgras{plaaslike pers} voer 'n veldtog
sodat daardie beligtingstelsel verbeter word, \VUgras{soos} 'n mens al \VUgras{die
trekstelsel van die} tremme volmaak het. \VUgras{Die} ou rytuie, \VUgras{getrek} deur
\VUgras{perde,} is prakties vervang deur \VUgras{elektriese tremme}
\textsuperscript{(31)} wat die reisigers vinnig van die een uiteinde van die \VUgras{stad
na die ander} bring, teen 'n baie \VUgras{goedkoop} prys.

%%K t11-c1-10-1 p
10. --- Naby die Geregshof is die \VUgras{Bank} \textsuperscript{(32)} waar 'n mens die
handelswaardes verdiskonteer. Die \VUgras{bankbeamptes} \textsuperscript{(33)} gaan die
skulde by die huise self in.

%%K t11-tit-4 sub
\VUsaut{4.0mm}
\VUpk{10.2pt}{40}{Kuns en onderwys.}
\VUsaut{3.4mm}

%%K t11-c1-11-1 p
11. --- Die mooi gebou wat daarnaas staan, is die \VUgras{museum.} Dit bevat tegelyk die
\VUgras{biblioteek} \textsuperscript{(34)} van die stad, die mooi
\VUgras{natuurwetenskaplike versamelings} \textsuperscript{(35)} (Natuurmuseum) en 'n
sierlike \VUgras{skildermuseum} \textsuperscript{(36)} wat 'n pragtige blywende
\VUgras{tentoonstelling} vorm.

%%K t11-c1-11-2 p
Dit gaan twee keer per week vir die publiek oop, maar die vreemdelinge mag dit enige dag
besoek teen 'n fooitjie aan die \VUgras{bewaker} \textsuperscript{(37)}. Die
\VUgras{skilders} \textsuperscript{(38)} kom daar werk. 'n Mens sien hulle \VUgras{voor}
hulle esel, met 'n \VUgras{palet} in die hand, besig om die beroemde skilderye af te
skryf.

%%K t11-c1-12-1 p
12. --- Op die hoogte, agter die museum, begin 'n mens die \VUgras{koepel}
\textsuperscript{(40)} van die \VUgras{hospitaal} \textsuperscript{(39)} te sien en die
geboue van die \VUgras{normaalskool} \textsuperscript{(41)} waarin die onderwysers van
ons gemeenteskole opgelei is. Op die Teaterplein is 'n \VUgras{poskantoor}
\textsuperscript{(43)} wat in 'n \VUgras{private huis} \textsuperscript{(42)} gevestig is.

%%K t11-tit-5 sub
\VUsaut{5.0mm}
\VUpk{11.4pt}{40}{Die brand.}
\VUsaut{3.0mm}

%%K t11-tit-6 sub
\VUpk{10.2pt}{40}{Brandkreet.}
\VUsaut{3.4mm}

%%K t11-c2-01-1 p
1. --- 'n Verskriklike gerug versprei hom deur die stad: daar het pas brand in die teater
uitgebreek! Op die oomblik dat ek op die \VUgras{plein} \textsuperscript{(96)} aankom, is
die skare al op die \VUgras{sypaadjie} \textsuperscript{(46)} en op die \VUgras{rypad}
\textsuperscript{(47)} versamel. Die \VUgras{posbeamptes} \textsuperscript{(44)} en die
\VUgras{briewebestellers} \textsuperscript{(45)} kom by die poskantoor uit. 'n
\VUgras{Trembestuurder} \textsuperscript{(48)} moes sy tremwa laat stilhou om 'n
stedelike \VUgras{ambulans} \textsuperscript{(50)} te laat verbygaan, wat met 'n
\VUgras{chirurg} \textsuperscript{(52)} en 'n \VUgras{verpleër} \textsuperscript{(51)}
aankom om 'n \VUgras{gewonde} \textsuperscript{(53)} te versorg wat op 'n
\VUgras{draagbaar} \textsuperscript{(54)} naby die \VUgras{koerantkiosk}
\textsuperscript{(55)} gebring is.

%%K t11-c2-02-1 p
2. --- 'n Kompanie infanterie, wat van die oefeninge af teruggekeer het, het stilgehou om
die ordediens te verseker saam met die \VUgras{berede stadspolisie}
\textsuperscript{(61)}. \VUgras{Die ruiters,} wie se perde steier, laat, beter as die \VUgras{soldate}
\textsuperscript{(57)} of die \VUgras{gendarmes} \textsuperscript{(63)}, die
\VUgras{nuuskieriges} \textsuperscript{(56)} agteruitgaan. Origens het die
gendarmes dit baie besig. Een van hulle wys die \VUgras{polisiemanne}
\textsuperscript{(64)} waarheen 'n mens 'n ander \VUgras{beseerde} \textsuperscript{(65)}
moet vervoer, 'n dapper brandweerman wat van 'n leer afgeval het.

%%K t11-c2-03-1 p
3. --- \VUgras{Die offisier} \textsuperscript{(66)} gee die plek noukeurig aan waar 'n
mens die \VUgras{stoomspuit} \textsuperscript{(67)} moet plaas wat aankom. Terwyl hy op
daardie kragtige masjien wag, het hy 'n \VUgras{handspuit} \textsuperscript{(68)} in
stelling gebring waarvan die lang \VUgras{slang} \textsuperscript{(69)} die water soos 'n
stortvloed op die vuur werp. Ses brandweermanne bedien dit, terwyl hulle maat, wat die
\VUgras{straalpyp} \textsuperscript{(70)} vashou, die \VUgras{straal}
\textsuperscript{(71)} na die middelpunt van die brand rig.

%%K t11-c2-04-1 p
4. --- Aan die ander kant van die teater het 'n mens die groot \VUgras{leer}
\textsuperscript{(72)} teen gesit. Dit maak dit vir die brandweermanne moontlik om die
vlamme te nader wat tussen warrelinge van swart \VUgras{rook} \textsuperscript{(75)}
opspuit.

%%K t11-tit-7 sub
\VUsaut{5.0mm}
\VUpk{10.2pt}{40}{Dapper dade.}
\VUsaut{3.4mm}

%%K t11-c2-05-1 p
5. --- \VUgras{Die brand} \textsuperscript{(74)} het in die voorportaal van die
\VUgras{teater} \textsuperscript{(73)} ontstaan, terwyl 'n mens \VUgras{die opera}
geoefen het waarvan die eerste \VUgras{opvoering} môre sou plaasgevind het. Gelukkig was
daar maar min \VUgras{toeskouers} \textsuperscript{(76)} in die saal. Die ongelukkiges
het na die \VUgras{balkon} \textsuperscript{(77)} gevlug waar dapper
\VUgras{brandweermanne} \textsuperscript{(86)} hulle met 'n leer en toue te hulp kom.

%%K t11-c2-06-1 p
6. --- Onder die \VUgras{peristiel} \textsuperscript{(78)}, waar die \VUgras{plakkaat}
\textsuperscript{(79)} oor die opvoering inlig, sien 'n mens die \VUgras{musikante}
\textsuperscript{(81)} van \VUgras{die orkes} \textsuperscript{(80)} vlug, terwyl hulle
hulle instrumente saamdra: \VUgras{viool} \textsuperscript{(81)}, \VUgras{fluit}
\textsuperscript{(82)}, \VUgras{tjello} \textsuperscript{(83)}, \VUgras{tromboon}
\textsuperscript{(84)}, \VUgras{trom} \textsuperscript{(85)}, ens.. 'n Mens probeer 'n
deel van die \VUgras{meubels} \textsuperscript{(87)} red.

%%K t11-c2-07-1 p
7. --- \VUgras{Die akteurs} \textsuperscript{(89)} en \VUgras{die aktrises}
\textsuperscript{(88)}, \VUgras{sangers} en \VUgras{sangeresse} moes in hulle \VUgras{toneelkostuum}
\textsuperscript{(90)} vlug. Die tenoor verduidelik aan 'n \VUgras{joernalis}
\textsuperscript{(92)} hoe dit gebeur het. \VUgras{Die verslaggewer} het dadelik van die
eerste oomblik af aangehardloop saam met die gesagsdraers: die \VUgras{prefek}
\textsuperscript{(91)}, die \VUgras{burgemeester} \textsuperscript{(93)}, die
\VUgras{polisiekommissaris} \textsuperscript{(94)} en die \VUgras{geregtelike amptenaar}
\textsuperscript{(95)}, wat ondersoek sal instel om oor die verantwoordelikhede te
oordeel.
\VUsaut{6.0mm}
\VUfilet{22mm}
